\documentclass[UTF8,a4paper,zihao=-4]{ctexart}

\usepackage{geometry}
\geometry{left=2.5cm,right=2.5cm,top=2.5cm,bottom=2.5cm}
\usepackage{booktabs}
\usepackage{longtable}
\usepackage{array}
\usepackage{enumitem}
\usepackage{hyperref}
\usepackage{caption}

\setlength{\tabcolsep}{4pt}

\hypersetup{
  colorlinks=true,
  linkcolor=black,
  urlcolor=blue
}

\setlist[itemize]{leftmargin=2em,itemsep=0.2em,topsep=0.3em}
\setlist[enumerate]{leftmargin=2em,itemsep=0.2em,topsep=0.3em}

\title{二手交易系统需求分析报告}
\author{需求分析工程师：241840078 张惠泽\\甲方：吴嘉聪}
\date{2026 年 10 月 18 日}

\begin{document}

\maketitle

\section{项目概况}

\subsection{要解决的问题}
\begin{itemize}
  \item 商品信息由卖家自行发布，质量与合规性缺少把关，需要管理员审核、下架与违规处理机制；
  \item 同一商品可能有多件库存，多个买家同时预约时容易出现重复预约、超卖与旧页面误操作，需要保证数据的一致性；
  \item 买卖双方在付款与发货之间存在信任缺口，在交易之间需要平台托管货款，确认收货后再向卖家打款；
  \item 交易达成后仍可能出现不满意或违规交易，需可追溯的沟通记录与退货退款通道；
\end{itemize}

因此，本系统要解决的问题是：在买卖双方之间建立一个规则明确、状态可追踪、并发安全、数据一致性能得到保证的线上二手交易流程，具体见下述项目目标：

\subsection{项目目标}

\begin{enumerate}
  \item 明确的角色规则：明确买家、卖家、管理员三类角色的权限边界。
  \item 覆盖商品从草稿、审核、发布、修改、下架到违规隐藏的完整生命周期。
  \item 科学的运行流程：以“预约—卖家确认—限时付款—平台托管—确认收货—打款”为主线，定义交易流程及其超时、取消与退货分支。
  \item 并发安全和多端数据一致性：对并发预约、重复预约、旧页面操作给出可执行的冲突处理规则与量化验收标准，同时多端数据要尽量快速地保证一致性。
\end{enumerate}

\subsection{系统范围}

本期系统包含以下功能范围：

\begin{itemize}
  \item 用户与权限：买家、卖家自助注册登录；管理员账号由系统预置；角色权限与互斥规则。
  \item 商品管理：草稿箱、提交审核、批准发布、修改重审、下架、违规隐藏、状态与可见性控制。
  \item 商品查询：按关键词、分类、价格区间、状态查询；按匹配度、近期评价、浏览量、销量加权排序。
  \item 交易流程：预约、取消预约、超时自动取消、卖家确认、限时付款、交易完成。
  \item 资金流程：平台托管货款、确认收货后打款、长期未确认自动打款、退货退款。
  \item 沟通与监管：买卖双方消息沟通、消息留存、管理员在举报场景下查看会话。
  \item 非功能约束：前后端分离与敏感数据保护、多端功能与数据一致、并发冲突处理与旧页面失效。
\end{itemize}

\subsection{明确不实现的内容}
\begin{itemize}
  \item 真实支付渠道、银行卡绑定与真实资金结算；支付采用模拟形式实现。
  \item 物流平台接口对接与自动签收；快递单号仅由买家录入，可由管理员人工确认。
  \item 买卖双方绕过平台的线下交易（不考虑）。
  \item 商品不设置延迟售卖。预约仅定义为发起购买意向。
\end{itemize}

\section{整体业务分析}

\subsection{角色与权限}

\begin{table}[htbp]
\centering
\caption{角色与主要权限}
\begin{tabular}{p{2.4cm}p{8.6cm}p{3.6cm}}
\toprule
角色 & 主要权限 & 限制 \\
\midrule
买家 & 浏览与查询商品、预约、取消预约、付款、确认收货、发起退货、与卖家沟通 & 不能创建或审核商品 \\
卖家 & 创建草稿、提交审核、修改商品、下架自有商品、确认交易、与买家沟通 & 不能审核商品，不能购买自己发布的商品 \\
管理员 & 审核商品、禁售/隐藏违规商品、查看被举报商品的会话、裁决退货争议 & 不参与买卖；不能与买家或卖家身份叠加 \\
\bottomrule
\end{tabular}
\end{table}

一个账号可以同时具备买家与卖家身份；管理员账号由系统预置，不开放自助注册，且不能与买家或卖家身份叠加。

\subsection{主要业务流程}

\subsubsection{商品从草稿到发布}

\begin{table}[htbp]
\centering
\caption{商品发布流程}
\begin{tabular}{p{1.2cm}p{2.4cm}p{5.2cm}p{5.4cm}}
\toprule
步骤 & 执行者 & 操作 & 结果 \\
\midrule
1 & 卖家 & 新建商品，信息不完整时保存草稿 & 商品处于草稿状态，进入草稿箱 \\
2 & 卖家 & 提交审核 & 商品进入待审核，买家不可见 \\
3 & 管理员 & 审核通过 & 商品发布，卖家与管理员收到通知 \\
4 & 管理员 & 审核驳回 & 商品退回草稿，驳回理由必填并通知卖家 \\
5 & 卖家 & 修改已发布商品 & 进入待审核；审核期间买家不可见，已有预约终止 \\
6 & 管理员 & 修改审核通过 & 更新线上内容并恢复可见 \\
7 & 管理员 & 修改审核驳回 & 回滚至修改前版本并恢复可见，通知卖家 \\
\bottomrule
\end{tabular}
\end{table}

草稿箱最多保存 10 篇草稿，草稿自最后一次保存起保留 30 天；超限时拒绝新建，到期前 3 天提醒，到期后自动清理。

\subsubsection{预约到交易完成}

\begin{table}[htbp]
\centering
\caption{预约与交易主流程}
\begin{tabular}{p{1.2cm}p{2.4cm}p{5.4cm}p{5.2cm}}
\toprule
步骤 & 执行者 & 操作 & 结果 \\
\midrule
1 & 买家 & 查询商品并查看详情与余量 & 展示符合条件的商品及综合排序 \\
2 & 买家 & 发起预约 & 服务端校验后余量减 1，生成预约记录 \\
3 & 买卖双方 & 沟通交易细节 & 消息送达，记录留存 30 天 \\
4 & 卖家 & 在预约后 24 小时内确认 & 交易进入待付款 \\
5 & 买家 & 在确认后 15 分钟内付款 & 交易完成，货款进入平台托管 \\
6 & 买家 & 确认收货 & 托管货款打给卖家，交易结束 \\
7 & 系统 & 买家 7 天未确认收货 & 自动打款给卖家 \\
\bottomrule
\end{tabular}
\end{table}

预约时服务端须校验：买家已登录、商品已发布、余量大于 0、商品非本人发布、该买家未重复预约同一商品。任一项不满足即拒绝并提示；页面展示的余量过期时，以服务端实际余量为准。

\subsubsection{取消与超时}

\begin{table}[htbp]
\centering
\caption{取消与超时分支}
\begin{tabular}{p{4.2cm}p{4.6cm}p{5.2cm}}
\toprule
触发条件 & 处理 & 结果 \\
\midrule
买家手动取消预约 & 释放预约占用的库存 & 余量加 1，通知买卖双方 \\
预约后 24 小时内买家未完成付款 & 系统自动取消 & 余量加 1，通知买卖双方 \\
预约后 24 小时内卖家未确认 & 系统自动取消 & 余量加 1，通知买卖双方 \\
卖家确认后 15 分钟内买家未付款 & 系统自动取消 & 余量加 1，通知买卖双方 \\
卖家确认后买家要求取消 & 不允许直接取消 & 只能走退货或等待付款超时 \\
\bottomrule
\end{tabular}
\end{table}

\subsubsection{下架、违规与退货}

卖家可下架自有商品；管理员可下架商品或判定违规并隐藏，隐藏与下架对买家的效果相同。商品下架或隐藏后，已有预约终止并通知买家、释放余量；已付款交易由管理员参考双方意见裁决：继续交易则按原流程执行，终止交易则进入退货退款流程。

\begin{table}[htbp]
\centering
\caption{退货流程}
\begin{tabular}{p{1.2cm}p{2.4cm}p{5.2cm}p{5.4cm}}
\toprule
步骤 & 执行者 & 操作 & 结果 \\
\midrule
1 & 买家 & 发起退货申请 & 进入待管理员审核 \\
2 & 管理员 & 审核退货申请 & 通过则待买家寄回；驳回则交易维持 \\
3 & 买家 & 寄回商品并填写快递单号 & 进入已寄回 \\
4 & 管理员 & 确认快递送达 & 等待卖家确认收货 \\
5 & 卖家 & 确认收到退货商品 & 平台将托管货款退还买家 \\
6 & 系统 & 买家未按期寄回 & 视为交易达成，不退货 \\
\bottomrule
\end{tabular}
\end{table}

卖家不同意退货时，由管理员根据双方理由裁决；买家未按期寄回的，“按期”的具体天数需在后续细化中明确。

\subsection{主要数据及其关系}

\begin{table}[htbp]
\centering
\caption{主要数据实体与关系}
\begin{tabular}{p{2.2cm}p{5.4cm}p{6.6cm}}
\toprule
实体 & 关键内容 & 主要关系 \\
\midrule
用户 & 账号、角色、联系方式 & 一个卖家对应多个商品；一个买家对应多个预约；买卖双方之间产生多条消息 \\
商品 & 名称、分类、价格、库存、状态、版本号 & 属于一个卖家；对应多条预约、消息与评分记录 \\
草稿 & 未发布商品的编辑内容，以商品状态标识 & 属于一个卖家，受草稿箱容量与保存期限约束 \\
预约 & 买家、商品、数量、预约时间、状态 & 多个预约对应一个商品；付款后生成订单 \\
订单 & 金额、付款时间、托管状态、收货状态 & 与预约一一对应；对应一条托管记录，可对应多条退货申请 \\
消息 & 发送者、接收者、商品、内容、时间 & 围绕一个商品，在买卖双方之间传递 \\
托管记录 & 应收、已付、待打款、已打款金额 & 与订单一一对应 \\
退货申请 & 理由、审核状态、快递单号 & 属于一个订单 \\
评分记录 & 评分值、评价时间 & 属于一个已完成订单，并汇总到商品排序 \\
\bottomrule
\end{tabular}
\end{table}

\subsection{业务状态与状态转换}

\subsubsection{商品状态}

\begin{table}[htbp]
\centering
\caption{商品状态转换}
\begin{tabular}{p{2.8cm}p{3.0cm}p{2.8cm}p{2.4cm}p{3.4cm}}
\toprule
当前状态 & 触发事件 & 执行者 & 目标状态 & 说明 \\
\midrule
草稿 & 提交审核 & 卖家 & 待审核 & 买家不可见 \\
待审核 & 审核通过 & 管理员 & 已发布 & 通知卖家与管理员 \\
待审核 & 审核驳回 & 管理员 & 审核驳回 & 理由必填并通知卖家 \\
审核驳回 & 修改后重新提交 & 卖家 & 待审核 & 每天提交不超过 5 次 \\
已发布 & 库存扣减至 0 & 系统 & 售罄 & 仍可被搜索，不可预约 \\
售罄 & 预约取消、余量恢复 & 系统 & 已发布 & 恢复可预约 \\
已发布/售罄 & 下架 & 卖家 & 已下架 & 买家不可见 \\
已发布/售罄 & 判定违规并隐藏 & 管理员 & 违规隐藏 & 仅管理员可恢复 \\
已下架 & 重新提交审核 & 卖家 & 待审核 & 审核期间不可见 \\
违规隐藏 & 复核通过 & 管理员 & 已发布 & 恢复对外可见 \\
已发布 & 修改商品 & 卖家 & 待审核 & 审核期间不可见，已有预约终止；驳回则回滚旧版本 \\
\bottomrule
\end{tabular}
\end{table}

\subsubsection{预约与订单状态}

\begin{table}[htbp]
\centering
\caption{预约与订单状态转换}
\begin{tabular}{p{2.9cm}p{4.2cm}p{4.2cm}p{3.2cm}}
\toprule
当前状态 & 触发事件 & 目标状态 & 说明 \\
\midrule
待卖家确认 & 卖家在 24 小时内确认 & 待付款 & 通知买家付款 \\
待卖家确认 & 买家取消，或 24 小时未付款 & 已取消 & 释放余量 \\
待卖家确认 & 卖家 24 小时未确认 & 已取消 & 释放余量 \\
待付款 & 买家在 15 分钟内付款 & 已付款（托管中） & 货款进入平台 \\
待付款 & 15 分钟未付款 & 已取消 & 释放余量 \\
已付款（托管中） & 买家确认收货 & 已完成 & 打款给卖家 \\
已付款（托管中） & 7 天未确认收货 & 已完成 & 系统自动打款 \\
已付款（托管中） & 退货流程完成 & 已退款 & 托管货款退还买家 \\
\bottomrule
\end{tabular}
\end{table}

\subsubsection{退货状态}

\begin{table}[htbp]
\centering
\caption{退货状态转换}
\begin{tabular}{p{2.8cm}p{4.4cm}p{3.2cm}p{4.0cm}}
\toprule
当前状态 & 触发事件 & 目标状态 & 说明 \\
\midrule
无 & 买家发起退货申请 & 待管理员审核 & 交易暂停结算 \\
待管理员审核 & 管理员通过 & 待买家寄回 & 通知买家填写快递单号 \\
待管理员审核 & 管理员驳回 & 无（交易维持） & 通知买卖双方 \\
待买家寄回 & 买家填写快递单号 & 已寄回 & 等待送达确认 \\
已寄回 & 管理员确认送达且卖家确认收货 & 已退款 & 托管货款退还买家 \\
待买家寄回/已寄回 & 超过期限未寄回 & 无（交易维持） & 视为交易达成 \\
\bottomrule
\end{tabular}
\end{table}

\subsection{并发与旧页面处理}

预约、取消、卖家确认、付款、草稿保存等写操作均可能发生并发冲突。系统采用服务端乐观锁与页面令牌：写操作携带资源版本号，版本不一致即拒绝并返回“版本冲突/页面已过期”，页面在 1 秒内提示并刷新到最新数据。库存与状态始终以服务端实际值为准，不允许旧页面覆盖新版本数据。验收以 100 个并发请求争抢 1 件库存为基准：成功数为 1、失败数为 99，重复 10 轮结果一致，超卖次数为 0。


\section{实验总结}

\subsection{访谈中的困难与解决方法}

访谈主要存在三类困难：

\begin{enumerate}
  \item 口语化与信息分散。受访者的表述包含重复、停顿和跳跃，同一规则可能分散在多个时间点。解决方法是先保留原始转写，再整理润色版，并为每段对话标注行号（如 INT-01-L211），使需求条目可以逐条回溯到原话。
  \item 前后表述不一致。例如“预约后 24 小时内未付款自动取消”与“卖家确认后 15 分钟内付款”在计时口径上并不自然衔接；乙方提出的“确认与交易完成是否只属于购买环节”也未在访谈中当场澄清。解决方法是将冲突显式记录为待确认项，通过补充沟通确认后再统一时间线。
  \item 边界情况缺失。草稿箱容量、草稿保存期限、审核不通过的处理、旧页面失效的判定方式等，访谈中只给出方向性描述。解决方法是先按同一字段结构登记为待确认需求，再逐项补充沟通，避免自行臆造规则。
\end{enumerate}

\subsection{需求整理中的困难与解决方法}

需求整理阶段的主要困难是“把自然语言变成可验收的条款”：

\begin{itemize}
  \item 非功能需求容易写成“安全”“稳定”“体验良好”。解决方法是将每条非功能需求拆成对象、条件与判断标准三个部分，并给出可执行的数量指标与测试手段。
  \item 商品状态在访谈中散落于草稿、审核、售罄、下架、违规等段落，缺少统一的状态全集。解决方法是建立状态转换表，明确每个状态对买家、卖家、管理员的可见性与可操作性。
  \item 初期出现了同一目标被拆成多条需求的情况，例如草稿箱规则、排序权重、旧页面失效实现方式各自单列。解决方法是以“一条需求一个主要目标”为准合并重复条目，删除被吸收的编号，同时保持其余编号不变。
  \item 异常分支写了但验收方法没有覆盖。解决方法是逐条对照“异常或边界情况”补充验收用例，使每个异常都有对应的可执行检查。
\end{itemize}

\subsection{人工审查的作用}

人工审查采用逐条复核的方式，每次记录相关需求、AI 原始内容、发现的问题、核对依据、最终处理与修改后内容。典型记录包括：

\begin{itemize}
  \item REV-000：消息形式与留存要求原为待确认，经补充沟通明确支持文字与图片、记录留存 30 天。
  \item REV-001：REQ-017 原缺少可测的正常结果，审查后补充对象、条件与四条量化判断标准。
  \item REV-002：补充沟通确认“预约即发起购买意向，先预约、卖家确认后付款”，据此删除重复的 REQ-024，并清理关联引用。
\end{itemize}

人工审查的价值在于判断“哪些内容是甲方真正说过的、哪些是整理者补的”，并保证需求状态、正文与关联引用三者一致。

\subsection{AI 辅助需求整理的作用与局限}

AI 在本项目中的作用主要体现在四个方面：

\begin{enumerate}
  \item 结构化：把口语化访谈快速整理为统一字段的需求条目，并保留来源行号。
  \item 一致性检查：发现同一规则在不同条目中的表述冲突，例如计时口径、管理员注册方式、重复定义的禁售职责。
  \item 可测化建议：为非功能需求生成对象、条件与判断标准，为异常分支生成验收用例。
  \item 维护性检查：检查失效引用、缺失字段、状态与正文不一致等文档问题。
\end{enumerate}

AI 的局限同样明显：

\begin{itemize}
  \item 会自行补全访谈中未说明的细节。例如草稿容量、保存期限、排序权重等数值属于建议值，必须经甲方确认后才能视为需求。
  \item 容易把“合理推测”写成“已确认”，导致状态与内容不一致；需要人工逐条核对来源并更新状态。
  \item 对业务语义的把握不稳定，可能保留自相矛盾的表述，或把实现方式误写成需求。
  \item 批量修改后容易产生失效引用与编号空缺，需要工具化检查（全文检索旧编号、统计状态与类型数量）。
\end{itemize}

因此，AI 适合承担整理、检查和起草工作，需求是否成立仍须由需求分析工程师与甲方共同确认。

\section{版本信息}

\begin{table}[htbp]
\centering
\caption{仓库与阶段提交}
\begin{tabular}{p{3.4cm}p{11.2cm}}
\toprule
项目 & 内容 \\
\midrule
仓库地址 & \url{https://github.com/zhangway123/homework_software_engineering} \\
lab1-ai-draft & \texttt{ab7ee2d0fb5ca88d63496ebd822109633dd4eeea} \\
lab1-final & \texttt{0cf38fd4fba5b563efe983e3be4ace4a96f88e7e} \\
\bottomrule
\end{tabular}
\end{table}

\end{document}
